\chapter{Stochastic Processes}

Many phenomena in nature, engineering, and economics evolve in time and are subject to
random influences. Typical examples include thermal motion of particles, fluctuating
signals, population dynamics, and financial time series. In such situations, randomness
cannot be captured by a single random variable; instead, one needs to describe an
entire \emph{family} of random quantities indexed by time or space.

A stochastic process provides a mathematical framework to model such systems. Informally,
it assigns to each time $t$ a random variable $X_t$ describing the state of the system
at that time. However, when dealing with infinitely many times, this informal description
raises fundamental questions: on which probability space do all these random variables
live, how are they related to each other, and what does it mean to specify their joint
distribution?

The approach taken in this section is to answer these questions in a systematic and
measure--theoretic way. Rather than postulating the existence of a stochastic process
from the outset, we start by describing its behavior at finitely many times. This leads
naturally to the notion of \emph{finite--dimensional distributions}, which encode all
probabilistic information about the process when observed at a finite number of time
indices.

The central issue is then \emph{consistency}: the distributions prescribed at different
levels of resolution must be compatible with one another. For example, the distribution
of $(X_{t_1},X_{t_2})$ must agree with the marginal distribution of
$(X_{t_1},X_{t_2},X_{t_3})$ when the third coordinate is ignored. Without such compatibility,
no single stochastic process can realize all prescribed distributions simultaneously.

The Daniell--Kolmogorov theorem shows that this intuitive requirement is not only necessary
but also sufficient. It provides a rigorous construction of a probability measure on the
space of all trajectories $E^T$ from a consistent family of finite--dimensional
distributions. In this way, the theorem forms the mathematical foundation for the
existence of stochastic processes and justifies the path--space viewpoint adopted in
this section.

\section{Basic definitions}

The definitions and results that follow should be understood as preparatory
steps toward the fundamental construction in Daniell-Kolgomorov theorem. Cylinder sets, consistency conditions, and
regularity assumptions on the state space are introduced precisely to make the
this theorem applicable and to ensure that stochastic processes can be
treated as genuine random elements in an infinite--dimensional setting.

\begin{definition}[Stochastic process]
Let $(\Omega,\mathcal{F},\Pp)$ be a probability space, let $(E,\B(E))$ be a measurable space,
and let $T$ be an index set.

A \emph{stochastic process} with state space $E$ and index set $T$ is a family
\[
(X_t)_{t\in T}
\]
of $E$--valued random variables on $(\Omega,\mathcal{F},\Pp)$, that is, measurable maps
\[
X_t : (\Omega,\mathcal{F}) \to (E,\B(E)) \qquad \text{for each } t\in T.
\]
\end{definition}

\begin{remark}[Path--space formulation of a stochastic process]
Let $(E,\mathcal{B}(E))$ be a measurable space and let $T$ be an index set.

We denote by
\[
E^T := \{ \omega : T \to E \}
\]
the set of all functions from $T$ into $E,$
called the \emph{path space}. An element $\omega \in E^T$ represents a single realization
(or trajectory) of the stochastic process.

The $\sigma$--algebra $\mathcal{B}(E^T)$ is defined as the smallest
$\sigma$--algebra such that for every $t \in T$ the coordinate projection
\[
\pi_t : E^T \to E, \qquad \pi_t(\omega) := \omega(t),
\]
is measurable.

With this notation, a stochastic process $(X_t)_{t \in T}$ can equivalently be described
as a single measurable mapping
\[
X : \Omega \to E^T,
\qquad
X(\omega) := (X_t(\omega))_{t \in T}.
\]

The equivalence follows from the fact that $X$ is measurable if and only if each
coordinate mapping $X_t = \pi_t \circ X$ is measurable.
\end{remark}


\subsection{Finite--Dimensional Distributions}

When studying a stochastic process indexed by a possibly infinite set $T$, it is in
general neither feasible nor necessary to describe the joint distribution of the entire
collection $(X_t)_{t\in T}$ at once. Instead, one starts by analyzing the behavior of the
process when observed at finitely many time points. The resulting joint laws encode all
probabilistic information that can be accessed through finite observations and are therefore
the most natural and fundamental building blocks of a stochastic process. Finite--dimensional
distributions allow us to describe local statistical behavior without committing to a
specific underlying probability space, and they provide a concrete and tractable way to
formulate existence questions. The central problem addressed later is whether a family of
such finite descriptions can be assembled into a single global object, a question that
leads naturally to the notions of marginal distributions, consistency, and ultimately to
the Daniell--Kolmogorov theorem.

But before approaching all these questions, we define Finite-dimensional distributions and consistency formally.

\begin{definition}[Finite--dimensional distributions]
Let $I = \{t_1,\dots,t_n\} \subset T$ be a finite subset.
The \emph{finite--dimensional distribution} of $(X_t)_{t\in T}$ at times $I$
is the probability measure
\[
\mu_I := \Pp \circ (X_{t_1},\dots,X_{t_n})^{-1}
\]
on $(E^I,\B(E^I))$ where $E^I$ is the product space defined by the index $I$.
\end{definition}

\begin{definition}[Consistency]
A family $(\mu_I)_{I\subset T,\ |I|<\infty}$ of probability measures on
$(E^I,\B(E^I))$ is called \emph{consistent} if for all finite $I \subset J \subset T$,
\[
\mu_I = \mu_J \circ \pi_{J\to I}^{-1},
\]
where $\pi_{J\to I} : E^J \to E^I$ denotes the canonical projection.
\end{definition}

\begin{remark}
Consistency expresses the requirement that lower--dimensional marginals are obtained
from higher--dimensional ones by projection.
\end{remark}

\paragraph{Consistency and Marginal Distributions.}

The consistency condition expresses the fact that all finite--dimensional distributions
describe different ``views'' of the \emph{same} stochastic process and therefore must be
compatible with each other.

To make this precise, let $I = \{t_1,\dots,t_n\} \subset J = \{t_1,\dots,t_n,t_{n+1},\dots,t_m\}$.
The probability measure $\mu_J$ describes the joint distribution of the random vector
\[
(X_{t_1},\dots,X_{t_n},X_{t_{n+1}},\dots,X_{t_m}).
\]

The corresponding distribution of $(X_{t_1},\dots,X_{t_n})$ alone is obtained by
\emph{marginalization}, that is, by integrating out the remaining variables:
\[
\mu_J \circ \pi_{J\to I}^{-1}(A)
=
\mu_J(A \times E^{J\setminus I}),
\qquad
A \in \mathcal{B}(E^I).
\]

\begin{remark}[Marginal distributions]
The measure $\mu_I$ is therefore interpreted as the \emph{marginal distribution} of $\mu_J$
with respect to the coordinates indexed by $I$.
\end{remark}

The family $(\mu_I)_{I \subset T}$ is called consistent precisely when the distribution
$\mu_I$ coincides with this marginal for every $I \subset J$, i.e.,
\[
\mu_I(A)
=
\mu_J(A \times E^{J\setminus I}),
\qquad
A \in \mathcal{B}(E^I).
\]

\paragraph{Why consistency is necessary?}
If a stochastic process $(X_t)_{t\in T}$ exists on some probability space, then its
finite--dimensional distributions must satisfy the consistency condition automatically.
Indeed,
\[
\Pp\big((X_t)_{t\in I} \in A\big)
=
\Pp\big((X_t)_{t\in J} \in A \times E^{J\setminus I}\big),
\]
which shows that $\mu_I$ is the marginal of $\mu_J$. In other words, observing more coordinates but imposing no restriction on them does not change the probability of an event that depends only on fewer coordinates. Consistency guarantees this almost obvious concept, i.e., if one adds more observations on time, then the probability of previous observations must remain unchanged.  

Consistency is therefore a \emph{necessary} condition for the existence of a stochastic
process with prescribed finite--dimensional distributions.

\paragraph{Why consistency is sufficient?}
The Daniell--Kolmogorov theorem shows that consistency is also \emph{sufficient}:
if the family $(\mu_I)$ is consistent, then there exists a probability measure on the
path space $E^T$ whose finite--dimensional distributions coincide with $(\mu_I)$.

In this sense, the consistency condition encodes exactly the information needed to
``glue together'' the finite--dimensional distributions into a single global object.

\paragraph{Intuitive interpretation.}
One may think of the measures $\mu_I$ as describing the behavior of a process when
observed at finitely many times. The consistency condition ensures that these observations
do not contradict each other when taken at different levels of resolution.

Without consistency, one could prescribe a distribution for $(X_1,X_2)$ and a different,
incompatible distribution for $X_1$ alone, making it impossible to construct a single
stochastic process realizing both.


\subsection{Cylinder Sets}

When constructing a stochastic process on the path space $E^T$, one faces an immediate
difficulty: the space $E^T$ is infinite--dimensional, and its Borel $\sigma$--algebra
cannot be described in a simple or concrete way. In particular, it is not feasible to
assign probabilities directly to arbitrary subsets of $E^T$.

The key observation is that all probabilistic information about a stochastic process is
obtained through finite observations. Events such as
\[
\{X_{t_1} \in B_1,\dots,X_{t_n} \in B_n\}
\]
depend only on the values of the process at finitely many time indices. Cylinder sets
are introduced precisely to formalize such events on the path space. They allow us to
encode finite--dimensional events as subsets of $E^T$ in a canonical way.

From a measure--theoretic point of view, cylinder sets play a dual role. On the one hand,
they form a concrete class of sets on which probabilities can be defined using the given
finite--dimensional distributions. On the other hand, they generate the product
$\sigma$--algebra on $E^T$, ensuring that a probability measure defined on cylinder sets
extends uniquely to all measurable events.

The Daniell--Kolmogorov theorem exploits this structure: a pre--measure is first defined
on an algebra of (special) cylinder sets using consistency of the finite--dimensional
distributions, and regularity properties of the state space then guarantee that this
pre--measure extends to a probability measure on the full path space. Cylinder sets thus
provide the essential link between finite--dimensional probabilistic descriptions and
the existence of a stochastic process as a random element of $E^T$.

\begin{definition}[Cylinder sets on the path space]
Let $(E,\mathcal{B}(E))$ be a measurable space and let $T$ be an index set.
Set
\[
\Omega := E^T.
\]

For a finite subset $I \subset T$ and a measurable set $A_I \in \mathcal{B}(E^I)$, we define
the \emph{cylinder set} associated with $A_I$ as
\[
C_{A_I}
:=
\bigl\{
\omega \in E^T \;:\; (\omega(t))_{t \in I} \in A
\bigr\}
= A_I \times E^{T \backslash I} = \pi_I^{-1}(A_I).
\]
\end{definition}


Cylinder sets describe events that have fixed finitely many coordinates of a path.

\subsection*{Special cylinder sets}

In order to work with a smaller and more manageable class of sets, we introduce
a special subclass of cylinder sets.

\begin{definition}[Special cylinder sets]
A \emph{special cylinder set} is a set of the form
\[
C
=
\bigcap_{k=1}^n \{ \omega \in E^T : \omega(t_k) \in B_k \} = \bigcap_{k=1}^n \pi_k^{-1}(B_k),
\]
where
\[
t_1,\dots,t_n \in T,
\qquad
B_1,\dots,B_n \in \mathcal{B}(E).
\]
\end{definition}

Equivalently, special cylinder sets are cylinder sets of the form
\[
C = \pi_I^{-1}(B_1 \times \cdots \times B_n) \times E^{T \backslash I},
\qquad
I = \{t_1,\dots,t_n\}.
\]

\begin{remark}
Special cylinder sets correspond to \emph{rectangular} sets in finite dimensions.
They are particularly convenient for constructing measures and checking additivity.
\end{remark}

\subsection*{Algebraic properties}

\begin{lemma}
The collection $\mathcal{A}$ of all special cylinder sets is an algebra of sets, i.e.,
\begin{enumerate}[label=(\roman*)]
\item $\Omega \in \mathcal{A}$,
\item $A \in \mathcal{A} \Rightarrow A^c \in \mathcal{A}$,
\item $A,B \in \mathcal{A} \Rightarrow A \cup B \in \mathcal{A}$.
\end{enumerate}
\end{lemma}

\begin{proof}
The whole space $\Omega$ corresponds to choosing $B_k = E$.
Complements and finite unions of special cylinder sets can be rewritten as finite unions
of special cylinder sets by enlarging the underlying finite index set and using
elementary set operations in $E$.
\end{proof}

\subsection*{Generation of the product $\sigma$--algebra}
A remarkable property of the $\sigma$--algebra generated by (special) cylinder
sets is that it coincides with the full product $\sigma$--algebra on $E^T$. Consequently, a
probability measure that is uniquely determined on cylinder sets extends uniquely to all
measurable subsets of the path space. This result provides the essential link between
finite--dimensional probabilistic information and the infinite--dimensional measurable
structure.
\begin{lemma}
The $\sigma$--algebra generated by the collection of special cylinder sets coincides with
the product $\sigma$--algebra:
\[
\sigma(\mathcal{A}) = \mathcal{B}(E^T).
\]
\end{lemma}

\begin{proof}
Every special cylinder set is a cylinder set, hence
\[
\sigma(\mathcal{A}) \subset \mathcal{B}(E^T).
\]

Conversely, the product $\sigma$--algebra $\mathcal{B}(E^T)$ is generated by sets of the
form $\pi_I^{-1}(A)$ with $A \in \mathcal{B}(E^I)$.
Since rectangles generate $\mathcal{B}(E^I)$ and rectangles correspond to special
cylinder sets, it follows that
\[
\mathcal{B}(E^T) \subset \sigma(\mathcal{A}).
\]
\end{proof}

\begin{remark}
The Daniell--Kolmogorov extension theorem is proved by first defining a pre--measure
on the algebra of special cylinder sets and then extending it uniquely to
$\mathcal{B}(E^T)$.
\end{remark}

\subsection{Inner Regularity and Tightness}
We will need a concept (or assumption) that control mass, when we pass limits in Daniell-Kolgomorov theorem. 
To construct a probability measure on the infinite--dimensional path space $E^T$ from
finite--dimensional distributions, it is not sufficient to work solely at the level of
algebras of sets. One must ensure that probabilities assigned to decreasing sequences of
events behave continuously, so that the resulting set function is $\sigma$--additive.
Inner regularity and tightness provide precisely the compactness properties needed to
control this behavior. They guarantee that probability mass cannot ``escape to infinity''
when passing to limits, and they allow events to be approximated from within by compact
sets with almost the same probability. In the context of the Daniell--Kolmogorov theorem,
these properties ensure that the pre--measure defined on cylinder sets extends to a
genuine probability measure on the full product $\sigma$--algebra. Inner regularity and
tightness therefore play a crucial technical role in turning finite--dimensional
consistency into a well--defined infinite--dimensional probability law.
\begin{lemma}[Inner regularity]
Let $E$ be a Polish space and let $\mu$ be a probability measure on $\B(E)$.
Then for every $A \in \B(E)$,
\[
\mu(A)
=
\sup\{\mu(K) : K \subset A,\; K \text{ compact}\}.
\]
\end{lemma}

\begin{remark}
In particular, every probability measure on a Polish space is tight.
This property is crucial to ensure $\sigma$--additivity when extending
pre--measures defined on algebras of sets.
\end{remark}

\section{Daniell--Kolmogorov Extension Theorem}

Up to this point, we have described a stochastic process only through its behavior at
finitely many time indices, encoded by its finite--dimensional distributions. This reflects
the way stochastic systems are observed in practice: one can measure or specify joint
distributions at finitely many times, but never at infinitely many times simultaneously.
The central question is whether such finite descriptions actually correspond to a genuine
stochastic process evolving over the entire index set $T$.

The difficulty lies in the infinite--dimensional nature of the path space $E^T$. While
finite--dimensional distributions live on well--understood product spaces, the space of
all trajectories is far more complex, and it is not a priori clear that one can define a
probability measure on it that is compatible with all finite observations. In particular,
different finite--dimensional distributions might contradict each other when viewed as
marginals of a single global distribution.

The Daniell--Kolmogorov theorem resolves this problem. It shows that the intuitive
compatibility requirement expressed by consistency of finite--dimensional distributions
is not only necessary but also sufficient for the existence of a stochastic process. If
the distributions prescribed at all finite sets of time indices agree with each other
under marginalization, then there exists a probability measure on the path space $E^T$
whose projections reproduce exactly these distributions.

In this sense, the Daniell--Kolmogorov theorem provides the mathematical justification for
the path--space approach to stochastic processes. It allows one to construct stochastic
processes starting solely from finite--dimensional probabilistic information, without
having to specify an underlying probability space in advance. All subsequent notions,
such as canonical processes, filtrations, and regularity properties of paths, rely on this
fundamental existence result.

\begin{theorem}[Daniell--Kolmogorov]
Let $E$ be a Polish space and let
\[
(\mu_I)_{I \subset T,\ |I|<\infty}
\]
be a consistent family of probability measures on $(E^I,\B(E^I))$.

Then there exists a unique probability measure
\[
\Pp \in \mathcal{P}(E^T,\B(E^T))
\]
such that for every finite $I \subset T$,
\[
\Pp \circ \pi_I^{-1} = \mu_I.
\]
\end{theorem}

\begin{proof}

\Step{1: Definition of a pre--measure}
For a cylinder set $C = \pi_I^{-1}(A)$ with $A \in \B(E^I)$, define
\[
\Pp_0(C) := \mu_I(A).
\]

\Step{2: Well--definedness}
If $\pi_I^{-1}(A) = \pi_J^{-1}(B)$, then both sets depend only on coordinates in $I\cup J$.
By consistency,
\[
\mu_I(A)
=
\mu_{I\cup J}(A \times E^{J\setminus I})
=
\mu_J(B),
\]
so $\Pp_0$ is well defined.

\Step{3: Finite additivity}
Let $(C_k)_{k=1}^n$ be disjoint cylinder sets.
By representing all $C_k$ on a common finite index set and using finite additivity of
$\mu_I$, we obtain
\[
\Pp_0\Big(\bigcup_{k=1}^n C_k\Big)
=
\sum_{k=1}^n \Pp_0(C_k).
\]

\Step{4: $\sigma$--additivity}
Since $E$ is Polish, the measures $\mu_I$ are inner regular and hence tight.
This prevents loss of mass at infinity and implies that $\Pp_0$ is
$\sigma$--additive on the algebra of cylinder sets.

\Step{5: Extension}
By the Carath\'eodory extension theorem, $\Pp_0$ extends uniquely to a probability
measure $\Pp$ on $\B(E^T)$.

\end{proof}

\subsection*{Existence of stochastic processes in Lusin spaces}

The Daniell--Kolmogorov theorem was stated under the assumption that the state space $E$
is Polish. This assumption can be relaxed to the broader class of Lusin spaces.

\begin{definition}[Lusin space]
A measurable space $(E,\mathcal{E})$ is called a \emph{Lusin space} if there exists
a Polish space $(\tilde E,\mathcal{B}(\tilde E))$ and a measurable bijection
\[
\varphi : E \to \tilde E
\]
such that both $\varphi$ and $\varphi^{-1}$ are measurable.
\end{definition}

\begin{remark}
Every Polish space is a Lusin space, but the converse is not true.
In a Lusin space, the measurable structure behaves essentially as well as in the Polish case,
even though the topology may be less regular.
\end{remark}

\begin{theorem}[Existence of stochastic processes in Lusin spaces]
Let $(E,\mathcal{E})$ be a Lusin space and let
\[
(\mu_I)_{I \subset T,\ |I|<\infty}
\]
be a consistent family of probability measures on $(E^I,\mathcal{E}^{\otimes I})$.

Then there exists a probability space $(\Omega,\mathcal{F},\Pp)$ and an $E$--valued
stochastic process $(X_t)_{t\in T}$ such that for every finite $I \subset T$,
\[
\Pp \circ (X_t)_{t\in I}^{-1} = \mu_I.
\]
\end{theorem}

\begin{proof}[Idea of the proof]
Since $(E,\mathcal{E})$ is a Lusin space, there exists a measurable isomorphism
$\varphi : E \to \tilde E$ onto a Polish space $\tilde E$.

For each finite $I \subset T$, define a probability measure $\tilde\mu_I$ on
$\tilde E^I$ by pushforward:
\[
\tilde\mu_I := \mu_I \circ \varphi_I^{-1},
\qquad
\varphi_I(x_1,\dots,x_n) := (\varphi(x_1),\dots,\varphi(x_n)).
\]

The family $(\tilde\mu_I)$ is consistent.
By the Daniell--Kolmogorov theorem in the Polish space $\tilde E$, there exists a probability
measure $\tilde\Pp$ on $\tilde E^T$ with finite--dimensional distributions $(\tilde\mu_I)$.

Finally, pulling back $\tilde\Pp$ via the inverse isomorphism $\varphi^{-1}$ yields
a probability measure $\Pp$ on $E^T$, and the canonical process
\[
X_t(\omega) := \omega(t)
\]
has finite--dimensional distributions $(\mu_I)$.
\end{proof}

\begin{remark}
The use of Lusin spaces is essential in applications where the state space is naturally
measurable but not Polish. The theorem shows that the existence of a stochastic process
depends primarily on the measurable structure rather than on the topology.
\end{remark}
